%=====================================================================
%  E-COMMERCE APPLICATION DEVELOPMENT --- FALL 2026
%  One-page weekly lecture schedule
%
%  SOURCE OF RECORD: parts/G_schedule.tex in the handout.  Change the
%  handout first, then this page -- the slide pipeline parses the
%  appendix, not this file.
%
%  Follows HEC 2017 (BS & MS (CS, SE & IT) Curriculum, p. 94), which is
%  the only published specification of this course in any HEC edition;
%  there it is an elective in BS Software Engineering.  The 2023 edition
%  drops it; the 2025 edition renames it "Digital Marketing &
%  E-Commerce" and publishes no contents.  The department teaches it as
%  INFT-4041, BSIT, 5th semester, 3 credit hours.
%
%  The preamble is shared with the sibling courses' schedule pages.
%
%  Build:  pdflatex ECAD_Fall2026_Lecture_Schedule.tex
%=====================================================================
\documentclass[10pt,a4paper]{article}
\usepackage[a4paper,left=1.2cm,right=1.2cm,top=0.7cm,bottom=0.45cm]{geometry}
\usepackage[T1]{fontenc}
\usepackage{lmodern}
\usepackage{microtype}
\usepackage[table]{xcolor}
\usepackage{array,ragged2e}
\usepackage{graphicx}
\usepackage{fontawesome5}
\pagestyle{empty}
\pdfinfo{/Title (E-Commerce Application Development, Fall 2026: Weekly Lecture Schedule)}
\setlength{\parindent}{0pt}
\setlength{\parskip}{0pt}

%--- palette: shared BSIT palette, part colours darkened a little so
%    white chip text stays legible at 6 pt
\definecolor{primary}{HTML}{1A5276}
\definecolor{secondary}{HTML}{2E86C1}
\definecolor{accent}{HTML}{C0392B}
\definecolor{lightbg}{HTML}{EBF5FB}
\definecolor{neutral}{HTML}{707B7C}
\definecolor{sub}{HTML}{465152}
\definecolor{p1}{HTML}{1A5276}
\definecolor{p2}{HTML}{2874A6}
\definecolor{p3}{HTML}{1E8449}
\definecolor{p4}{HTML}{8E44AD}
\definecolor{p5}{HTML}{117A65}
\definecolor{p6}{HTML}{8A6D0A}

%--- building blocks -------------------------------------------------
\newcommand{\chip}[2]{{\setlength{\fboxsep}{1.3pt}\colorbox{#1}{\makebox[1.45em][c]{%
  \fontsize{6.4}{7}\selectfont\bfseries\color{white}#2}}}}
\newcommand{\chipw}[2]{{\setlength{\fboxsep}{1.3pt}\colorbox{#1}{%
  \fontsize{6.4}{7}\selectfont\bfseries\color{white}\,#2\,}}}
% \ch{part}{number}{title}{what it covers}  -- one chapter, own paragraph
\newcommand{\ch}[4]{\chip{p#1}{#2}\,\textbf{#3}\ {\color{neutral}--}\ {\fontsize{6.9}{8.0}\selectfont\color{sub}#4}\par}
\newcommand{\ex}[1]{{\bfseries\color{accent}#1}}      % exams and the big deliverables
\newcommand{\ev}[1]{{\bfseries\color{primary}#1}}     % quizzes, assignments, project points
\newcommand{\none}{{\color{neutral}---}}
\newcommand{\hdr}[1]{{\bfseries\color{white}#1}}
\newcommand{\lab}[1]{{\color{sub}#1}}

\newlength{\topicw}   % set inside the document, once \linewidth and \tabcolsep are final
\newcolumntype{K}{>{\centering\arraybackslash}p{0.75cm}}
\newcolumntype{T}{>{\RaggedRight\arraybackslash}p{\topicw}}
\newcolumntype{L}[1]{>{\RaggedRight\arraybackslash}p{#1}}
% \leavevmode: a cell that opens with \color puts a whatsit before its paragraph, and
% \vtop then aligns it by its top edge -- one strut lower than a cell that opens with text
\newcommand{\wk}[4]{\leavevmode{\large\bfseries\color{primary}#1} & \leavevmode #2 &
  \leavevmode #3 & \leavevmode #4 \\ \hline}

\newcommand{\fact}[3]{{\setlength{\fboxsep}{4pt}\colorbox{lightbg}{%
  \parbox[t]{\dimexpr#1-8pt\relax}{\RaggedRight\fontsize{6.2}{7}\selectfont\bfseries\color{neutral}\MakeUppercase{#2}\\[1.5pt]%
  \fontsize{7.8}{9.2}\selectfont\bfseries\color{primary}#3}}}}
\newcommand{\tile}[3]{{\colorbox{lightbg}{\parbox[c][0.9cm][c]{\dimexpr(\linewidth-3pt*\numexpr#1-1\relax)/#1-6pt-0.5pt\relax}{%
  \centering{\large\bfseries\color{primary}#2}\\[-0.5pt]{\fontsize{6.6}{7.4}\selectfont\color{sub}#3}}}}}

% \legend{I}{II}{III}{IV}{V}{VI} -- the six part names, as a 3 x 2 grid of colour chips
\newcommand{\legend}[6]{{\fontsize{6.6}{8}\selectfont\color{sub}%
  \begin{tabular}{@{}p{0.31\linewidth}p{0.31\linewidth}p{0.31\linewidth}@{}}
  \chipw{p1}{I}\ #1 & \chipw{p2}{II}\ #2 & \chipw{p3}{III}\ #3 \\[1pt]
  \chipw{p4}{IV}\ #4 & \chipw{p5}{V}\ #5 & \chipw{p6}{VI}\ #6
  \end{tabular}}}

\begin{document}
\fontsize{7.9}{9.3}\selectfont

%--- header ----------------------------------------------------------
\begin{minipage}[c]{1.5cm}%
  \IfFileExists{assets/iub-logo.png}{\includegraphics[height=1.15cm]{assets/iub-logo.png}}{}%
\end{minipage}\hfill
\begin{minipage}[c]{0.65\linewidth}\centering
  {\footnotesize\bfseries\color{primary}THE ISLAMIA UNIVERSITY OF BAHAWALPUR}\\[1pt]
  {\scriptsize\color{neutral}Department of Information Technology}
\end{minipage}\hfill
\begin{minipage}[c]{1.5cm}\raggedleft
  \IfFileExists{assets/dit-logo.png}{\includegraphics[height=1.15cm]{assets/dit-logo.png}}{}%
\end{minipage}

\vspace{3pt}
{\color{secondary}\rule{\linewidth}{1.4pt}}\\[2pt]
% This course's name is fourteen characters longer than any sibling's, so
% the two halves of this line do not both fit at the family's usual sizes.
% The title keeps its weight and loses a step of size rather than wrapping.
{\large\bfseries\color{primary}E-COMMERCE APPLICATION DEVELOPMENT}\hfill
{\small\bfseries\color{secondary}Weekly Lecture Schedule\ \textbullet\ Fall 2026}\\[1pt]
{\footnotesize\color{sub}Bachelor of Science in Information Technology\ \textbullet\
 \emph{Building, Running and Securing an Online Shop}}\\[1pt]
{\footnotesize\color{sub}INFT-4041, 5th semester. One shop is installed in Week 3 and every week afterwards adds a layer to it; by Week 16 it is the project.}\\[1pt]
{\color{secondary}\rule{\linewidth}{0.6pt}}

\vspace{4pt}
\fact{0.235\linewidth}{Credit hours}{3 (3--0): 3 lecture hours a week}\hfill
\fact{0.235\linewidth}{Duration}{16 teaching weeks; midterm Week 8, project due Week 16}\hfill
\fact{0.235\linewidth}{Prerequisite}{Web Technologies (HEC names it Web Engineering)}\hfill
\fact{0.235\linewidth}{Contents}{23 chapters in 6 parts; 12 measured experiments on a live shop}

\vspace{5pt}
%--- the schedule ----------------------------------------------------
\begingroup
\renewcommand{\arraystretch}{0.96}
\setlength{\tabcolsep}{3.5pt}
\setlength{\topicw}{\dimexpr\linewidth-0.75cm-3.55cm-2.85cm-6\tabcolsep\relax}
\arrayrulecolor{primary!18}
\rowcolors{2}{white}{lightbg}
\noindent\begin{tabular}{@{}K T L{3.55cm} L{2.85cm}@{}}
\rowcolor{primary}
\hdr{\faIcon{calendar-alt}} & \hdr{\faIcon{book-open}\ \ Lecture topics} &
\hdr{\faIcon{flask}\ \ Laboratory} & \hdr{\faIcon{clipboard-check}\ \ Assessment} \\
\wk{1}
 {\ch{1}{1}{What E-Commerce Is, and What It Changed}{what changed and what did not; the four functions a shop performs}
  \ch{1}{2}{Business Models and Revenue Models}{B2C, B2B, C2C, marketplace; who holds the stock, the customer and the risk}}
 {\lab{Read four real shops; the arithmetic of one order}}{\none}
\wk{2}
 {\ch{1}{3}{Auctions, Portals and Communities}{the models HEC strands at the end of its list; what each one monetises}
  \ch{1}{4}{Planning an E-Commerce Framework}{requirements; build, buy or extend --- the decision the course rests on}}
 {\lab{Three shapes compared; your requirements and your decision}}{\none}
\wk{3}
 {\ch{2}{5}{The Stack, and Your Shop in One Command}{the four layers; containers, images and volumes; the shop is live this week}}
 {\lab{Compose up; verify, break and reset the shop}}{\none}
\wk{4}
 {\ch{2}{6}{Products, Categories and Their Tables}{the screens and the eighteen tables beneath them; what a catalogue costs as it grows}}
 {\lab{50{,}000 products, and the index that does not rescue the page}}{\none}
\wk{5}
 {\ch{2}{7}{Product Variations, Options and Uploads}{options grow by addition, variants by multiplication; uploads as a security topic}}
 {\lab{Option combinations, and what an upload field accepts}}{\ev{Quiz 1} (Ch 1--5)}
\wk{6}
 {\ch{2}{8}{Themes, Layout and the Customer's Experience}{a theme is a directory; override, never edit; what a page weighs, and why it is not images}}
 {\lab{Page weight before and after, in bytes and requests}}{\ev{Assignment 1 set}}
\wk{7}
 {\ch{3}{9}{The Shopping Basket, and Where State Lives}{cookie, session or table; what a session costs; what survives a restart}}
 {\lab{Session, cookie and database baskets across a restart}}{\none}
\wk{8}
 {\ch{3}{10}{The Checkout and the Order State Machine}{what each step commits; where the order row is written; status 0}}
 {\lab{Midterm, then queries counted per checkout step}}{\ex{Midterm} (Ch 1--9)}
\wk{9}
 {\ch{3}{11}{Shipping, Tax and Zones}{geo zones; rate, class and rule; the tax base; Pakistani sales tax}
  \ch{3}{12}{Discounts, Vouchers and Referrals}{four devices; totals are a pipeline, and position decides the bill}}
 {\lab{Geo zones and GST; a coupon applied in the wrong order}}{\ev{Assignment 1 due}}
\wk{10}
 {\ch{4}{13}{Web Payment Systems}{the five parties; authorisation, capture, settlement; what COD really costs}
  \ch{4}{14}{Taking Payment for Orders}{the handshake; a timeout from a distribution; idempotency; PCI DSS}}
 {\lab{A sandbox gateway, and what happens when it times out}}{\ev{Project proposal}}
\wk{11}
 {\ch{4}{15}{User Account Management}{why a password is stored slowly; cost factor against logins per second; what a login leaks}}
 {\lab{Password cost against logins per second}}{\none}
\wk{12}
 {\ch{4}{16}{Securing the Shop}{the legal boundary first; fifteen checks before and after one config file}}
 {\lab{The OWASP checks against your own container, before and after}}{\ev{Quiz 2} (Ch 10--15)}
\wk{13}
 {\ch{5}{17}{Administration: The Dashboard}{what each widget counts and why they disagree; the figure that changes a decision}
  \ch{5}{18}{Managing Orders, Customers and Refunds}{stock moves on transitions, not statuses; the two things called a refund}}
 {\lab{A refund through the order state machine}}{\ev{Assignment 2 set}}
\wk{14}
 {\ch{5}{19}{Extending the Platform: Your First Module}{configure, extend, edit --- in that order; a module and a shipping method, and what each costs}}
 {\lab{A shipping method and a dashboard widget, from nothing}}{\none}
\wk{15}
 {\ch{5}{20}{Deployment, Maintenance and Performance}{throughput against latency; the saturation point; what a restore costs}
  \ch{6}{21}{SEO and Digital Marketing}{what is checkable; readable URLs are two halves; what a crawler is told}}
 {\lab{Load test with the cache on and off; crawlability measured}}{\ev{Assignment 2 due}}
\wk{16}
 {\ch{6}{22}{Social, Legal and Ethical Issues}{what makes an order a contract; the Pakistani instruments; dark patterns}
  \ch{6}{23}{The Semester Project}{deploy from nothing, measure, present, defend; review another team}}
 {\lab{Project presentations}}{\ex{Project due; Final} (all)}
\end{tabular}
\endgroup

\vspace{4pt}
%--- assessment weights ----------------------------------------------
{\fontsize{6.6}{8}\selectfont\bfseries\color{neutral}\MakeUppercase{Assessment weights}}\\[2pt]
\tile{4}{30\%}{Sessional: 2 quizzes, 2 assignments, defence}\hfill
\tile{4}{25\%}{Semester project, built from Week 3}\hfill
\tile{4}{20\%}{Mid-term, Week 8 (Ch 1--9)}\hfill
\tile{4}{25\%}{Final-term, all chapters}

\vspace{3pt}
%--- legend and notes ------------------------------------------------
\legend{E-Commerce and Its Models}{Standing the Shop Up}{From Basket to Order}
       {Money and Identity}{Running and Extending the Shop}{Reach, Law and Practice}

\vspace{1pt}
{\fontsize{6.3}{7.6}\selectfont\color{neutral}%
Follows HEC 2017 (\emph{BS \& MS (CS, SE \& IT) Curriculum}, p.\ 94), the only published specification of this course in any HEC edition, where it is an \emph{elective in BS Software Engineering}; the 2023 edition drops the course entirely and the 2025 edition renames it \emph{Digital Marketing \& E-Commerce} with no contents. The department teaches it as INFT-4041, BSIT, 5th semester. All 22 distinct clauses of the 2017 outline are traced to chapters in the handout's Course Specification, which also records that the outline's ``Course Content'' paragraph is the table of contents of its own reference book~2 (Peacock, \emph{PHP 5 E-commerce Development}, 2010), duplicated clauses and all, and that HEC publishes only two CLOs, with the Bloom level column left blank. Three topics are taught that the outline does not name, each justified there: the data model behind the catalogue (Week 4), security as its own chapter (Week 12), and module development (Week 14), which is what the department's added words ``Application Development'' denote. HEC's prerequisite is ``Web Engineering''; the department's equivalent is Web Technologies at 4th semester. Chapter numbers refer to the handout, \emph{E-Commerce Application Development --- Building, Running and Securing an Online Shop} (219 pp., 23 chapters, 7 appendices); Appendix~G of that handout is the source of record for this page. Twelve chapters carry a measured experiment whose figures were produced on the machine described in Appendix~A and pasted into the text; the sources are in \texttt{code/} and each is reproducible. Every security exercise is confined to the student's own container on \texttt{localhost}; testing any other system is an offence under PECA 2016.\par}

\end{document}
